% UNINTEGRATED ARITHMETIC15 FOUNDATION PART, 25 October2026.
% Same existing arithmetic paper family; source preparation only.
% First batch: accepted general dihedral3/19 theorem and D14 consumer.
% DRAFT pending independent and coordinating review; no execution/compilation.
\section{A strengthened ordinary dihedral comparison}
\label{app:arithdihedral}

This section supplies the ordinary rank-two comparison needed beside
\(D_{14}\).  Its component theorem is general real geometry.
Arithmeticity enters only when the already established arithmetic
\(D_{14}\) prefix bound is used in the final corollary.
Throughout, \(p\) is the ORIGINAL vertex farthest from one strict
interior point.  The theorem is not asserted at an arbitrary
neighbour with a weakened comparison factor.

\begin{proposition}\label{prop:ordinarydihedralthreenineteenths}
Let a compact hyperbolic Coxeter polytope have an actual farthest
vertex \(p\), and write \(z=E(p-x)\), \(\|x\|<1\).
Let one WHOLE component of the incident finite diagram be
\(I_2(m)\), \(m\ge4\), and let \(e\) be the squared norm of the
projection of this SAME \(x\) onto that component.
If \(e\le3/19\), then \(m=7\) and one actual opposite wall
closes both component directions, has a single old contact of
label three there, and has no contact with any other old component.
This is precisely the actual \(O_7\) profile.
Consequently every ordinary \(I_2(m)\), \(m\ge4\), has energy
strictly greater than \(3/19\), with arbitrary finite spectators.
\end{proposition}

\begin{proof}
Use outward unit incident roots, finite Gram \(B\), inverse
\(C=B^{-1}\ge0\), duals \(\omega_i\), and unit inward dual
directions \(-u_i=-\omega_i/\sqrt{C_{ii}}\).
For an actual outgoing edge of length \(L_i\), maximality at
the original \(p\) gives
\[
 x\cdot u_i\ge\tau_i=\tanh(L_i/2)>0,\qquad
 r_i=\tanh L_i={2\tau_i\over1+\tau_i^2}.
\]
The old dihedral inverse is
\[
 C_I=d\begin{pmatrix}1&k\\k&1\end{pmatrix},\qquad
 k=\cos(\pi/m)\ge1/\sqrt2,\quad d=(1-k^2)^{-1}.
\]
Thus each \(\tau_i^2\le e\), and the increasing function
\(4t/(1+t)^2\) on \(0\le t<1\) gives
\[
                         r_i^2\le57/121=:r .
\tag{Dih1}\label{eq:arithdihedralr}
\]

We recall the actual-wall calculation forcing a common opposite.
For an outside actual wall \(c=Ap-v\), let
\(\beta_j=-c\cdot n_j\ge0\) and put
\[
 T_i(c)={(C\beta)_i\over A\sqrt{C_{ii}}}.
\]
For the actual opposite \(c_i\) of edge \(i\),
\(T_i(c_i)=1/r_i\).
The matrix
\(C-C_{\cdot j}C_{j\cdot}/C_{jj}\)
is entrywise nonnegative: its nonzero block is the inverse of
the positive principal Gram obtained by deleting \(j\).
Actual acuteness between two distinct opposites therefore gives
\(T_j(c_i)\le r_j\), while that same nonnegative matrix gives
\(T_j(c_i)\ge k/r_i\) within this old component.
Hence distinct opposites imply \(r_ir_j\ge k\).
But~\eqref{eq:arithdihedralr} gives
\(r_ir_j\le57/121<1/\sqrt2\le k\).
One actual wall \(c\) therefore closes BOTH directions.

Write its full decomposition as
\[
 c=Ap-v_I-v_S,\quad h=A^2>0,\quad
 b=-c\cdot n_1\ge0,\quad g=-c\cdot n_2\ge0,\quad
 Q=\|v_S\|^2\ge0 .
\]
Here \(S\) is the entire orthogonal spectator forest.
Every positive contact is a native magnitude at least \(1/2\);
none of its finite, parallel or dotted possibilities is omitted
before the actual endpoint constraints are imposed.
At least one of \(b,g\) is positive.

Suppose both are positive.
At the actual endpoint releasing direction one, the retained
white Gram has residual \(\delta=1-Q-g^2>0\).
Its actual edge formula is
\[
 \cosh L_1={\sqrt d(b+kg)\over\sqrt{1-Q-g^2}}
 \ge{\sqrt d(1+k)\over\sqrt3}
 ={\,\cot(\pi/(2m))\over\sqrt3}
 \ge{1+\sqrt2\over\sqrt3}>{11\over8}.
\]
The last comparison is \(128\sqrt2>171\), whose positive
squared margin is \(32768-29241\).
It follows that
\[
 e\ge\tanh^2(L_1/2)>
 {11/8-1\over11/8+1}={3\over19},
\]
a contradiction.  Hence exactly one old contact \(\beta>0\)
remains.

We need one whole-spectator fact at small time.
If \(h<1/4\), then \(c\) cannot contact any foreign component.
Indeed a positive foreign contact gives a normalized derivative
\[
 {(C_S\beta)_j\over\sqrt{(C_S)_{jj}}}\ge1/2 .
\]
This uses the full spectator inverse, with nonnegative entries
and diagonal at least one.
Its ratio to \(\sqrt h\) exceeds one.
The actual-wall shadow inequality therefore forces this SAME
\(c\) to be the opposite of that foreign edge as well.

For clarity, the required orthogonal obstruction is specific to
the original farthest vertex.  Suppose a wall closed orthogonal
dual directions \(u_i,u_j\).
Clear it positively off all other incident roots \(W\):
\[
 c_F=c+\sum_{s\in W}\alpha_sn_s,\qquad
 \alpha=B_W^{-1}(-c\cdot n_W)\ge0.
\]
Strict interiority gives \(z\cdot c_F<0\).
In the retained Lorentz three-space its exact form is
\[
 c_F=A p-(A/r_i)u_i-(A/r_j)u_j .
\]
Therefore
\[
 {x\cdot u_i\over r_i}+{x\cdot u_j\over r_j}<1.
\]
Original maximality gives the opposite bound
\(\tau_i/r_i+\tau_j/r_j
 =1+(\tau_i^2+\tau_j^2)/2>1\).
This proves the obstruction, including arbitrary contacts in \(W\).
Thus \(h<1/4\) really forces \(Q=0\).

Put \(q=d\beta^2\).
The smaller normalized derivative square is \(k^2q\).
The common-wall edge formula and~\eqref{eq:arithdihedralr} imply
\[
 q-1\le h=q+Q-1\le rk^2q .
\]
Substituting \(k^2=1-1/d\) and \(r=57/121\) gives the exact
strengthened restriction
\[
                         \beta^2(64d+57)\le121 .
\tag{Dih2}\label{eq:arithdihedralcut}
\]
This is derived at the \(3/19\) threshold; it is not the older
\(1/7\) inequality with an unchanged coefficient.

Since \(\beta^2\ge1/4\), this forces \(d\le427/64\).
For every old \(m\ge8\),
\(d\ge4+2\sqrt2>34/5>427/64\), impossible.
The complete remaining old alphabet is \(4,5,6,7\).
At old \(m=5\),
\(d=2+2/\sqrt5>185/64\):
equivalently \(128>57\sqrt5\), with squared margin139.
Since \(d\) increases with \(m\), if
\(\beta\ge1/\sqrt2\) at any of old \(5,6,7\),
\eqref{eq:arithdihedralcut} would give the opposite
\(d\le185/64\).
Thus all three rows have exactly \(\beta=1/2\).
Every larger finite contact, parallel and dotted contact has
been excluded by the native jump, rather than a sampled label list.

At old four, \(d=2\) and
\(\beta^2\le121/185<\cos^2(\pi/5)\).
The strict comparison is \(413<185\sqrt5\), with squared
margin556.  Only contacts three and four remain, giving
\(q=1/2\) or \(1\).
Consequently \(0<h\le57/242<1/4\).
The small-time argument just proved gives \(Q=0\), whereas
\(h=q-1\le0\), a contradiction.

At old five only \(\beta=1/2\) remains.
Since \(d<3\), \(q<3/4\) and
\(0<h\le r(d-1)/4<57/242<1/4\).
Again \(Q=0\), and \(h=q-1<0\).

At old six, \(\beta=1/2\), \(d=4\), \(q=1\).
If \(Q=0\), then \(h=0\), impossible.
For any \(Q>0\), the actual positive sum
\[
                          \nu=c+2n_1+\sqrt3n_2
\]
is a nonzero future null: all three selected products vanish.
The coefficient order is for OLD6/NEW3, since
\(c\cdot n_1=-1/2\), \(n_1\cdot n_2=-\sqrt3/2\),
\(c\cdot n_2=0\).
Every other actual wall pairs nonpositively termwise with \(\nu\).
It is a feasible ideal ray, contrary to compactness.
No foreign contact was discarded in this affine exclusion.

At old seven only \(\beta=1/2\) remains, and
\[
 0\le Q=1+h-d/4
 \le{427-64d\over484}<{107\over484}<{1\over4},
\]
using \(d_7>5\).
Every positive spectator cost is at least \(1/4\), by inverse
positivity and the native contact lower bound.  Thus \(Q=0\).
We have exactly the ACTUAL isolated-new-three wall in old
\(I_2(7)\), closing both directions: the stated \(O_7\).
Its time square is \(d_7/4-1\in(0,1/3)\);
the bounds \(5<d_7<16/3\) follow from
\(6/5<2\cos(2\pi/7)<5/4\).
Its separate inherited energy floor \(>9/100\) is unchanged.

All weak inequalities included \(e=3/19\), so no ordinary equality
case is omitted.  The proof has retained every finite spectator
component, all its contacts and the same original \(x\).
No arithmetic hypothesis or weakened-neighbour comparison was used.
\end{proof}

\begin{corollary}\label{cor:D14ordinarydihedral}
At the original farthest vertex in arithmetic dimension sixteen,
a whole \(D_{14}\) cannot have an ordinary irreducible
rank-two spectator.
\end{corollary}
\begin{proof}
The fresh true-rank-\(14\) argument in the proof of
Proposition~\ref{prop:f16} gives
\[
                         e_D>.84816={5301\over6250}
\]
for its actual whole rank-two spectator.
Its bootstrap uses only the independently available spectator
floor \(>.09\), so it does not depend on the present strengthening.
For an ordinary \(I_2(m)\), \(m\ge4\), the proposition gives
\[
 \|x\|^2=e_D+e_S>
 {5301\over6250}+{3\over19}
 =1+{719\over118750}>1 .
\]
For ordinary \(A_2\), use its geometric companion bound
\(319/2000\), which exceeds \(3/19\) by \(61/38000\).
Both contradict \(\|x\|<1\).
The actual typed \(N_7/O_7\) alternatives retain their separate
costs and arguments; they have not been charged as ordinary.
\end{proof}

% Next batches may append reviewed arithmetic proof bodies here.
% Preserve this exact first-batch source before any such expansion.

